\section{基准算例设计与评价协议}
为评估所提出方法在几何表达能力、力学预测能力与数值稳定性方面的性能，本文采用由基线模型、标准算例和统一评价指标组成的验证协议。完成稿应确保所有方法在相同几何、材料参数、时间积分设置和接触触发准则下进行比较，以避免由于实现细节差异而引入偏置。

\subsection{对比基线}
为验证本文方法的有效性，计划设置以下基线：
\begin{enumerate}[label=(\arabic*)]
    \item 单点罚函数接触模型；
    \item 多点离散罚函数模型；
    \item 传统点接触互补模型或约束模型；
    \item 必要时增加简化压力场模型作为中间对照。
\end{enumerate}

\subsection{典型验证场景}
本文计划选取以下场景进行测试：
\begin{enumerate}[label=(\arabic*)]
    \item 立方体平面接触与偏心受压；
    \item 立方体落角并翻转接触；
    \item 非凸刚体滚过边缘或凹槽；
    \item 薄壁零件接触；
    \item 双复杂 CAD 刚体大面积接触。
\end{enumerate}

\subsection{评价指标}
评价指标拟包括：最大压入量、净接触力误差、净接触力矩误差、接触响应连续性与高频抖动程度、对接触斑面积和接触中心位置的估计能力，以及时间步敏感性与计算开销。

\subsection{公平比较与结果报告原则}
为保证比较结论具有可解释性，完成稿应统一报告以下设置：
\begin{enumerate}[label=(\arabic*)]
    \item 几何来源、网格分辨率与 SDF 分辨率；
    \item 时间积分器、时间步长及收敛/截断阈值；
    \item 接触刚度、阻尼和摩擦参数的标定方式；
    \item 采样密度、自适应策略和接触历史滤波设置；
    \item 计算平台、单步耗时和总仿真耗时统计口径。
\end{enumerate}
